% !TEX program = pdflatex
\documentclass[11pt,a4paper]{article}

% =========================================================================
% ARQUITECTURA TÉCNICA
% Proyecto: CasaScore (AlquilaPuno)
% Gestión de Proyectos TICS (EST315) -- Semestre 2026-II -- Team Sobrevivientes
% =========================================================================

% --- PAQUETES DE CONFIGURACIÓN ---
\usepackage[utf8]{inputenc}
\usepackage[T1]{fontenc}
\usepackage[spanish,es-tabla]{babel}
% Compatibilidad con babel-spanish reciente (flechas de TikZ y signo de porcentaje)
\AtBeginDocument{\shorthandoff{<>}\spanishplainpercent}
\usepackage{geometry}
\geometry{
    top=3.6cm,
    bottom=2.3cm,
    left=2.2cm,
    right=2.2cm,
    % PARTE MODIFICABLE: headheight debe ser mayor que la altura del logo más alto (2.56cm = 73pt).
    headheight=80pt,
    headsep=15pt,
    footskip=20pt
}

\usepackage{graphicx}
\usepackage{fancyhdr}
\usepackage[table]{xcolor}
\usepackage{titlesec}
\usepackage{booktabs}
\usepackage{tabularx}
\usepackage{array}
\usepackage{amsmath}
\usepackage{enumitem}
\usepackage{microtype}
\usepackage{amssymb}
% --- RECURSOS VISUALES: GRÁFICOS Y CÓDIGO DE COLORES ---
\usepackage{tikz}
\usetikzlibrary{positioning,calc}
\definecolor{GrisFila}{RGB}{246, 248, 251}
\definecolor{CuadCerca}{RGB}{150, 30, 30}
\definecolor{CuadSatisf}{RGB}{180, 130, 20}
\definecolor{CuadInfo}{RGB}{10, 45, 90}
\definecolor{CuadMonit}{RGB}{110, 115, 125}
\newcommand{\etiqueta}[2]{{\setlength{\fboxsep}{2pt}\colorbox{#1}{\textcolor{white}{\scriptsize\bfseries #2}}}}
\newcommand{\cerca}{\etiqueta{CuadCerca}{Gestionar de cerca}}
\newcommand{\satisf}{\etiqueta{CuadSatisf}{Mantener satisfecho}}
\newcommand{\informado}{\etiqueta{CuadInfo}{Mantener informado}}
\newcommand{\monit}{\etiqueta{CuadMonit}{Monitorear}}
\usepackage{needspace}
\usepackage{longtable}
\definecolor{Rama1}{RGB}{10, 45, 90}
\definecolor{Rama2}{RGB}{40, 95, 150}
\definecolor{Rama3}{RGB}{180, 130, 20}
\definecolor{Rama4}{RGB}{30, 110, 60}
\definecolor{Rama5}{RGB}{110, 60, 140}
\definecolor{RaciR}{RGB}{150, 30, 30}
\definecolor{RaciA}{RGB}{10, 45, 90}
\definecolor{RaciC}{RGB}{180, 130, 20}
\definecolor{RaciI}{RGB}{140, 145, 155}
\definecolor{NivAlto}{RGB}{30, 110, 60}
\definecolor{NivMedio}{RGB}{180, 130, 20}
\definecolor{NivBajo}{RGB}{150, 30, 30}
\usetikzlibrary{patterns}
\definecolor{SevBaja}{RGB}{30, 110, 60}
\definecolor{SevMedia}{RGB}{180, 130, 20}
\definecolor{SevAlta}{RGB}{150, 30, 30}
\usetikzlibrary{arrows.meta,shapes.geometric,fit,backgrounds}
\usepackage{listings}
\lstset{basicstyle=\ttfamily\scriptsize, frame=single, rulecolor=\color{AzulUNA!40}, backgroundcolor=\color{GrisFila}, columns=fullflexible, keepspaces=true, breaklines=true, xleftmargin=4pt, xrightmargin=4pt, literate={á}{{\'a}}1 {é}{{\'e}}1 {í}{{\'i}}1 {ó}{{\'o}}1 {ú}{{\'u}}1 {ñ}{{\~n}}1 {Á}{{\'A}}1}

\usepackage{hyperref}

% Logos: funciona a 2 o 3 niveles de profundidad dentro de docs/latex/
\graphicspath{{../../assets/img/}{../../../assets/img/}}

% --- PALETA INSTITUCIONAL FIEI / UNA PUNO ---
\definecolor{AzulUNA}{RGB}{10, 45, 90}
\definecolor{DoradoFIEI}{RGB}{180, 130, 20}
\definecolor{GrisOscuro}{RGB}{50, 50, 50}
\definecolor{AzulClaro}{RGB}{228, 236, 246}

% --- CONFIGURACIÓN DE LISTAS ---
\setlist{topsep=2pt, itemsep=2pt, parsep=1pt}

% --- COLUMNAS DE TABLA ---
\newcolumntype{L}{>{\raggedright\arraybackslash}X}
\newcolumntype{P}[1]{>{\raggedright\arraybackslash}p{#1}}
\newcolumntype{C}[1]{>{\centering\arraybackslash}p{#1}}
\newcolumntype{Q}[1]{>{\raggedleft\arraybackslash}p{#1}}
\newcolumntype{R}{>{\raggedleft\arraybackslash}X}
\renewcommand{\arraystretch}{1.25}

% --- FORMATO DE ENLACES ---
\hypersetup{
    colorlinks=true,
    linkcolor=AzulUNA,
    urlcolor=AzulUNA,
    citecolor=AzulUNA,
    pdftitle={Arquitectura Técnica -- CasaScore (AlquilaPuno)},
    pdfauthor={Team Sobrevivientes}
}

% --- FORMATO DE TÍTULOS ---
\titleformat{\section}{\color{AzulUNA}\normalfont\large\bfseries}{\thesection.}{0.7em}{}[\color{AzulUNA}\titlerule]
\titleformat{\subsection}{\color{AzulUNA}\normalfont\normalsize\bfseries}{\thesubsection.}{0.5em}{}
\titlespacing*{\section}{0pt}{9pt}{4pt}
\titlespacing*{\subsection}{0pt}{6pt}{2pt}

% --- ENCABEZADO Y PIE DE PÁGINA (fancyhdr) ---
\pagestyle{fancy}
\fancyhf{}

% ===== TAMAÑO DE LOS LOGOS (PARTE MODIFICABLE) =====
% logo_fiei.png  : height=1.45cm (lado izquierdo)
% logo_equipo.png: height=2.56cm (lado derecho; 25 % más grande porque la imagen tiene más borde blanco)
% Para cambiar el tamaño, modifique solo el valor height=... de cada \includegraphics.
% Si el logo más alto supera 2.6cm, aumente también headheight y top en la configuración de página.
\fancyhead[L]{%
    \includegraphics[height=1.45cm,keepaspectratio]{logo_fiei.png}%
}
\fancyhead[C]{%
    \scriptsize \textbf{\textcolor{AzulUNA}{UNIVERSIDAD NACIONAL DEL ALTIPLANO DE PUNO}}\\[2pt]
    \tiny \textcolor{GrisOscuro}{FACULTAD DE INGENIERÍA ESTADÍSTICA E INFORMÁTICA}\\[2pt]
    \scriptsize \textbf{\textcolor{AzulUNA}{Gestión de Proyectos TICS (EST315) \textbar\ Team Sobrevivientes}}\\[2pt]
    \footnotesize \textbf{\textcolor{DoradoFIEI}{ARQUITECTURA TÉCNICA}}%
}
\fancyhead[R]{%
    \includegraphics[height=2.56cm,keepaspectratio]{logo_equipo.png}%
}

\renewcommand{\headrulewidth}{0pt}
\renewcommand{\headrule}{}

\fancyfoot[L]{\scriptsize \textcolor{GrisOscuro}{CasaScore (AlquilaPuno)}}
\fancyfoot[C]{\small \thepage}
\fancyfoot[R]{\scriptsize \textcolor{GrisOscuro}{Puno, Perú -- Semestre 2026-II}}
\renewcommand{\footrulewidth}{0.4pt}

\begin{document}

\begin{center}
    {\Large \textbf{\textcolor{AzulUNA}{ARQUITECTURA TÉCNICA}}}\\[3pt]
    {\normalsize \textbf{CasaScore (AlquilaPuno)}}\\[4pt]
    \footnotesize \textbf{Fecha:} martes 22/09/2026 (semana 5) \quad \textbar\quad \textbf{Versión:} 1.0 \quad \textbar\quad \textbf{Fase:} Planificación
\end{center}

\vspace{0.1cm}

\section{Datos Generales}
\noindent\begin{tabularx}{\textwidth}{@{}P{4.2cm} L@{}}
\toprule
\textbf{Proyecto:} & CasaScore (aplicación AlquilaPuno) \\
\textbf{Responsable:} & Eliseo Tarqui Ajahuana -- Especialista Técnico / Arquitectura \\
\textbf{Proceso PM4DEV:} & Gestión del alcance (diseño de la solución) \\
\textbf{Insumos:} & Especificación de requerimientos, plan de calidad y plan de contratos \\
\bottomrule
\end{tabularx}

\section{Visión General}
CasaScore es una aplicación móvil para Android que se conecta a un servicio en la nube. La arquitectura sigue tres decisiones:
\begin{itemize}[leftmargin=*]
    \item \textbf{Cliente con tres capas:} presentación (pantallas), lógica (motor de puntaje y distancia) y datos (acceso a la nube y caché local). Cada capa solo conoce a la de abajo.
    \item \textbf{Backend como servicio:} Supabase da cuentas, base de datos y almacenamiento de fotos, sin servidor propio.
    \item \textbf{Motor en el celular:} la app descarga los cuartos activos y calcula el CASA SCORE con las prioridades y el punto de referencia de cada arrendatario. El ranking responde rápido y no genera costo de servidor.
\end{itemize}

\section{Diagrama de Arquitectura}
\begin{center}
\begin{tikzpicture}[font=\scriptsize, >=Stealth, mod/.style={draw=#1, fill=#1!10, rounded corners=2pt, text width=2.55cm, align=center, minimum height=0.95cm, line width=0.6pt}]
\node[draw=AzulUNA, fill=AzulClaro!40, rounded corners=4pt, minimum width=9.6cm, minimum height=7.2cm] at (0,-0.3) {};
\node[font=\scriptsize\bfseries, text=AzulUNA] at (0,3.05) {APLICACIÓN ANDROID (Expo, React Native y TypeScript)};
\node[font=\tiny\bfseries, text=GrisOscuro, anchor=west] at (-4.65,2.45) {PRESENTACIÓN};
\node[mod=Rama2] at (-2.95,1.75) {Registro e inicio de sesión};
\node[mod=Rama2] at (0,1.75) {Ranking y detalle};
\node[mod=Rama2] at (2.95,1.75) {Publicar, perfil y mis cuartos};
\node[font=\tiny\bfseries, text=GrisOscuro, anchor=west] at (-4.65,0.95) {LÓGICA};
\node[mod=Rama3] at (-2.95,0.25) {\textbf{Motor CASA SCORE}\\ {\tiny scoring.ts}};
\node[mod=Rama3] at (0,0.25) {Distancia y tiempo\\ {\tiny geo.ts}};
\node[mod=Rama3] at (2.95,0.25) {Validación del formulario};
\node[font=\tiny\bfseries, text=GrisOscuro, anchor=west] at (-4.65,-0.55) {DATOS};
\node[mod=Rama4] (cli) at (1.5,-1.25) {Cliente de Supabase\\ {\tiny supabase.ts}};
\node[mod=Rama4] (cac) at (-1.5,-1.25) {Caché local SQLite\\ {\tiny database.ts}};
\draw[<->, Rama4] (cac) -- (cli);
\node[mod=Rama2] (map) at (1.5,-2.55) {Mapa\\ {\tiny MapLibre React Native}};
\node[draw=Rama5, fill=Rama5!8, rounded corners=4pt, text width=3.6cm, align=center, minimum height=3.2cm, line width=0.8pt] (sb) at (6.85,0.9) {\textbf{\textcolor{Rama5}{SUPABASE}}\\[3pt] Auth: cuentas y roles\\[2pt] PostgreSQL con seguridad por fila\\[2pt] Storage: fotos de los cuartos};
\node[draw=GrisOscuro, fill=GrisFila, rounded corners=4pt, text width=3.6cm, align=center, minimum height=1.4cm, line width=0.8pt] (osm) at (6.85,-2.55) {\textbf{OPENSTREETMAP}\\[2pt] Teselas del mapa de Puno};
\draw[<->, thick, Rama5] (cli.east) -- node[above, font=\tiny]{HTTPS} ++(1.55,0) |- (sb.west);
\draw[->, thick, GrisOscuro] (map.east) -- node[above, font=\tiny]{HTTPS} (osm.west |- map.east);
\end{tikzpicture}
\end{center}

\Needspace{14\baselineskip}
\section{Tecnologías}
{\small
\rowcolors{2}{white}{GrisFila}
\noindent\begin{tabularx}{\textwidth}{@{}P{3.6cm} L C{2.3cm}@{}}
\toprule
\rowcolor{AzulClaro}
\textbf{Tecnología} & \textbf{Uso} & \textbf{Estado} \\
\midrule
Expo SDK 57 y React Native 0.86 & Base de la aplicación Android & \etiqueta{Rama4}{En uso} \\
React 19.2 y TypeScript 6.0 & Interfaz y tipado estricto del código & \etiqueta{Rama4}{En uso} \\
expo-sqlite & Base local del MVP de la Unidad I y caché sin conexión & \etiqueta{Rama4}{En uso} \\
expo-image-picker & Captura de las fotos del cuarto & \etiqueta{Rama4}{En uso} \\
supabase-js & Conexión con Auth, base de datos y Storage & \etiqueta{SevMedia}{Semana 9} \\
MapLibre React Native & Mapa con teselas de OpenStreetMap, sin clave de Google & \etiqueta{SevMedia}{Semana 11} \\
expo-location & Ubicación actual al publicar un cuarto & \etiqueta{SevMedia}{Semana 11} \\
Jest & Pruebas automáticas del motor de puntaje & \etiqueta{SevMedia}{Semana 7} \\
\bottomrule
\end{tabularx}
}

\section{Módulos del Código}
{\small
\rowcolors{2}{white}{GrisFila}
\noindent\begin{tabularx}{\textwidth}{@{}P{4.2cm} L C{2.3cm}@{}}
\toprule
\rowcolor{AzulClaro}
\textbf{Archivo o carpeta} & \textbf{Responsabilidad} & \textbf{Estado} \\
\midrule
\texttt{src/screens/} & Pantallas: registro, ranking, detalle, publicar cuarto y perfil. & \etiqueta{SevMedia}{En avance} \\
\texttt{src/utils/scoring.ts} & Normalización, suma ponderada, bonificaciones, penalizaciones y desglose. & \etiqueta{SevMedia}{Por actualizar} \\
\texttt{src/utils/geo.ts} & Distancia y tiempo a pie al punto de referencia. & \etiqueta{SevAlta}{Por crear} \\
\texttt{src/data/database.ts} & Tabla local de cuartos en SQLite y datos de prueba. & \etiqueta{Rama4}{Terminado} \\
\texttt{src/data/supabase.ts} & Consultas a Supabase y subida de fotos. & \etiqueta{SevAlta}{Por crear} \\
\texttt{src/types/index.ts} & Tipos de cuarto, preferencias y desglose del puntaje. & \etiqueta{SevMedia}{Por actualizar} \\
\bottomrule
\end{tabularx}
}

\vspace{0.1cm}
\noindent El motor actual usa la distancia fija a la UNA y cuatro pesos. Se actualiza para leer el punto de referencia del arrendatario y agregar el peso del tipo de baño, según la especificación de requerimientos.

\Needspace{20\baselineskip}
\section{Modelo Entidad--Relación}
\begin{center}
\begin{tikzpicture}[>=Stealth]
\node[draw=AzulUNA, fill=AzulUNA!10, line width=0.8pt, text width=3.6cm, align=left, inner sep=3pt, font=\tiny] (usuario) at (0.00,0.00) {\textbf{\scriptsize\textcolor{AzulUNA}{USUARIO}}\\[-1pt]\rule{3.6cm}{0.3pt}\\ \underline{id} (uuid, de Supabase Auth)\\ nombre (texto)\\ correo (texto)\\ rol (arrendatario \textbar\ propietario)\\ creado\_en (fecha)};
\node[draw=Rama4, fill=Rama4!10, line width=0.8pt, text width=3.6cm, align=left, inner sep=3pt, font=\tiny] (preferencia) at (0.00,-3.60) {\textbf{\scriptsize\textcolor{Rama4}{PREFERENCIA}}\\[-1pt]\rule{3.6cm}{0.3pt}\\ \underline{usuario\_id} (FK)\\ peso\_precio, peso\_cercania\\ peso\_seguridad, peso\_servicios\\ peso\_bano (0 a 5)\\ ref\_latitud, ref\_longitud};
\node[draw=Rama3, fill=Rama3!10, line width=0.8pt, text width=3.9cm, align=left, inner sep=3pt, font=\tiny] (cuarto) at (5.60,0.00) {\textbf{\scriptsize\textcolor{Rama3}{CUARTO}}\\[-1pt]\rule{3.9cm}{0.3pt}\\ \underline{id} (uuid)\\ propietario\_id (FK a usuario)\\ titulo, direccion\\ precio (S/), area\_m2\\ bano\_privado, luz\_natural\\ latitud, longitud\\ estado (activo \textbar\ oculto)};
\node[draw=Rama5, fill=Rama5!10, line width=0.8pt, text width=3.2cm, align=left, inner sep=3pt, font=\tiny] (foto) at (11.00,0.90) {\textbf{\scriptsize\textcolor{Rama5}{FOTO}}\\[-1pt]\rule{3.2cm}{0.3pt}\\ \underline{id} (uuid)\\ cuarto\_id (FK)\\ tipo (fachada, habitacion,\\ bano, ventana, seguridad)\\ url (Supabase Storage)};
\node[draw=Rama2, fill=Rama2!10, line width=0.8pt, text width=3.2cm, align=left, inner sep=3pt, font=\tiny] (servicios) at (11.00,-2.20) {\textbf{\scriptsize\textcolor{Rama2}{SERVICIOS}}\\[-1pt]\rule{3.2cm}{0.3pt}\\ \underline{cuarto\_id} (FK)\\ agua\_24h, luz\_24h\\ internet\_24h};
\node[draw=SevAlta, fill=SevAlta!10, line width=0.8pt, text width=3.9cm, align=left, inner sep=3pt, font=\tiny] (seguridad) at (5.60,-3.90) {\textbf{\scriptsize\textcolor{SevAlta}{SEGURIDAD}}\\[-1pt]\rule{3.9cm}{0.3pt}\\ \underline{cuarto\_id} (FK)\\ porton\_llave, rejas\\ iluminacion, camara};
\node[draw=SevMedia, fill=SevMedia!10, line width=0.8pt, text width=3.2cm, align=left, inner sep=3pt, font=\tiny] (condiciones) at (11.00,-4.60) {\textbf{\scriptsize\textcolor{SevMedia}{CONDICIONES}}\\[-1pt]\rule{3.2cm}{0.3pt}\\ \underline{cuarto\_id} (FK)\\ entrada\_libre\_24h, hora\_llegada\\ visitas, cocina, mascotas\\ declaracion\_transparente};
\draw[thick, GrisOscuro] (usuario.east) -- node[above, font=\tiny, pos=0.45]{1 \ publica \ N} (cuarto.west);
\draw[thick, GrisOscuro] (usuario.south) -- node[right, font=\tiny]{1 : 1} (preferencia.north);
\draw[thick, GrisOscuro] (cuarto.east) -- node[above, font=\tiny, sloped]{1 : 4 o más} (foto.west);
\draw[thick, GrisOscuro] (cuarto.east) -- node[below, font=\tiny, sloped]{1 : 1} (servicios.west);
\draw[thick, GrisOscuro] (cuarto.south) -- node[right, font=\tiny]{1 : 1} (seguridad.north);
\draw[thick, GrisOscuro] (cuarto.south east) -- node[below, font=\tiny, sloped]{1 : 1} (condiciones.west);
\end{tikzpicture}
\end{center}

\Needspace{12\baselineskip}
\subsection{Seguridad por fila}
{\small
\rowcolors{2}{white}{GrisFila}
\noindent\begin{tabularx}{\textwidth}{@{}P{3.4cm} L L@{}}
\toprule
\rowcolor{AzulClaro}
\textbf{Tabla} & \textbf{Quién puede leer} & \textbf{Quién puede escribir} \\
\midrule
usuario & El propio usuario & El propio usuario \\
preferencia & El arrendatario dueño de la fila & El arrendatario dueño de la fila \\
cuarto y sus tablas & Todos los usuarios, solo cuartos activos & El propietario del cuarto \\
foto (Storage) & Todos los usuarios & El propietario del cuarto \\
\bottomrule
\end{tabularx}
}

\Needspace{16\baselineskip}
\section{Pantallas y Navegación}
\begin{center}
\begin{tikzpicture}[font=\scriptsize, >=Stealth, pa/.style={draw=#1, fill=#1!10, rounded corners=3pt, text width=2.6cm, align=center, minimum height=0.95cm, line width=0.7pt}]
\node[pa=AzulUNA] (reg) at (0,0) {\textbf{Registro e inicio de sesión}};
\node[pa=Rama4] (rk) at (4.2,1.1) {\textbf{Ranking}\\ {\tiny cuartos ordenados}};
\node[pa=Rama4] (dt) at (8.4,1.1) {\textbf{Detalle}\\ {\tiny 4 fotos y desglose}};
\node[pa=Rama4] (pf) at (12.3,1.1) {\textbf{Perfil}\\ {\tiny prioridades y punto}};
\node[pa=Rama3] (mc) at (4.2,-1.1) {\textbf{Mis cuartos}\\ {\tiny editar u ocultar}};
\node[pa=Rama3] (pb) at (8.4,-1.1) {\textbf{Publicar cuarto}\\ {\tiny formulario en 5 pasos}};
\node[pa=Rama3] (pp) at (12.3,-1.1) {\textbf{Perfil}\\ {\tiny datos de la cuenta}};
\draw[->, thick, Rama4] (reg) -- node[above, sloped, font=\tiny]{arrendatario} (rk);
\draw[->, thick, Rama3] (reg) -- node[below, sloped, font=\tiny]{propietario} (mc);
\draw[<->, Rama4] (rk) -- (dt); \draw[<->, Rama4] (dt) -- (pf);
\draw[<->, Rama3] (mc) -- (pb); \draw[<->, Rama3] (pb) -- (pp);
\end{tikzpicture}
\end{center}
\noindent Los 5 pasos de Publicar cuarto son: datos básicos; servicios y seguridad; condiciones de convivencia; ubicación en el mapa; y 4 fotos obligatorias. Sin un dato obligatorio no se puede avanzar al siguiente paso.

\section{Algoritmo Principal}
\subsection{Distancia y tiempo al punto de referencia}
La distancia en línea recta entre el cuarto $(\varphi_1,\lambda_1)$ y el punto de referencia $(\varphi_2,\lambda_2)$ se calcula con la fórmula de Haversine, con el radio de la Tierra $R = 6\,371\,000$ m:
\[
d = 2R \arcsin\!\sqrt{\sin^2\!\tfrac{\varphi_2-\varphi_1}{2} + \cos\varphi_1 \cos\varphi_2 \sin^2\!\tfrac{\lambda_2-\lambda_1}{2}}
\qquad\qquad
t = \frac{1.3\,d}{80}
\]
El tiempo a pie $t$ (minutos) considera que el recorrido por las calles es 30\,\% más largo que la línea recta y una velocidad de 80 m por minuto.

\Needspace{16\baselineskip}
\subsection{Cálculo del ranking}
\begin{lstlisting}
entrada: cuartos activos, prioridades w (5 pesos), punto de referencia
para cada cuarto c:
    d, t  <- distancia y tiempo de c al punto de referencia
    N     <- [precio(c), cercania(d, t), seguridad(c), servicios(c), bano(c)]   # 0 a 100
    base  <- suma(w[i] * N[i]) / suma(w)
    B     <- bonificaciones de c        # luz +8, entrada libre +5, 16 m2 o mas +4
    P     <- penalizaciones de c        # llegada, visitas, cocina, no declaradas
    score <- limitar(base + B - P, 0, 100)
    guardar desglose (N, B, P) para la pantalla de detalle
ordenar cuartos por score de mayor a menor
\end{lstlisting}
\noindent Cada cuarto se procesa una vez y luego se ordena la lista; con 500 cuartos el cálculo toma pocos milisegundos, dentro del límite de 250 ms.

\Needspace{14\baselineskip}
\section{Decisiones de Arquitectura}
{\small
\rowcolors{2}{white}{GrisFila}
\noindent\begin{tabularx}{\textwidth}{@{}C{1.3cm} P{3.6cm} L@{}}
\toprule
\rowcolor{AzulClaro}
\textbf{ID} & \textbf{Decisión} & \textbf{Motivo y consecuencia} \\
\midrule
DA-01 & Calcular el puntaje en el celular & Cada arrendatario tiene prioridades y un punto distinto; evita costo de servidor. A cambio, la app descarga los cuartos activos. \\
DA-02 & Supabase como backend & Cuentas, base de datos con seguridad por fila y fotos en un solo servicio gratuito. Depende de los límites de su plan. \\
DA-03 & OpenStreetMap con MapLibre & Mapa sin costo y sin tarjeta. Requiere una compilación de desarrollo de Expo en lugar de Expo Go. \\
DA-04 & SQLite en la Unidad I y Supabase en el MVP final & Permitió validar el motor sin conexión; luego SQLite queda como caché local. \\
DA-05 & Motor separado de las pantallas & Permite probarlo con Jest y cambiar las reglas sin tocar la interfaz. \\
\bottomrule
\end{tabularx}
}

\section{Límites del Software}
\begin{itemize}[leftmargin=*]
    \item Solo Android 10 o superior y solo la ciudad de Puno.
    \item La distancia es en línea recta y el tiempo a pie es una estimación; no se calculan rutas reales por las calles.
    \item El ranking está diseñado para hasta 500 cuartos activos.
    \item Se necesita conexión para iniciar sesión, publicar y actualizar los cuartos; sin conexión se muestra la última lista guardada.
    \item No verifica en persona los cuartos, no gestiona pagos ni contratos de alquiler y no garantiza la seguridad física del inmueble.
\end{itemize}

\Needspace{24\baselineskip}
\section{Firmas}
\noindent En señal de conformidad con el presente documento, los integrantes del equipo firman a continuación:

\vspace{1.5cm}
\begin{center}
\begin{tabular}{c c c}
    \parbox{6.5cm}{\centering \rule{6.0cm}{0.6pt}\\[5pt] \textbf{Clyde Neil Paricahua Pari}\\[2pt] {\footnotesize Project Manager}}
    & \hspace{0.8cm} &
    \parbox{6.5cm}{\centering \rule{6.0cm}{0.6pt}\\[5pt] \textbf{Ruth Karina Apaza Solis}\\[2pt] {\footnotesize Analista de Planificación}}
    \\[2.0cm]
    \parbox{6.5cm}{\centering \rule{6.0cm}{0.6pt}\\[5pt] \textbf{Eliseo Tarqui Ajahuana}\\[2pt] {\footnotesize Especialista Técnico / Arquitectura}}
    & \hspace{0.8cm} &
    \parbox{6.5cm}{\centering \rule{6.0cm}{0.6pt}\\[5pt] \textbf{Viviana Mary Quisbert Quispe}\\[2pt] {\footnotesize Gestora de Comunicación, Riesgos y Documentación}}
\end{tabular}
\end{center}
\end{document}
